\chapter{The Jagged Frontier}
\label{ch:jagged-frontier}

\epigraph{The person who knows where the frontier is jagged in their domain is the most valuable person in the organisation.}{}

\section{Why Wholesale Automation Is Reckless}

If coordination is now computationally cheap, the logical next question is whether to simply replace the hierarchy with AI. Several prominent voices --- most notably \textcite{DorseyBotha2026} --- propose exactly this. But the evidence reveals that the boundary between what AI can do and what it cannot is dangerously unpredictable. This boundary is the single most important concept in the entire coordination collapse thesis.

\section{The BCG Experiment}

The \meshterm{jagged technological frontier} was identified and named by \textcite{DellAcqua2026} in what remains the most rigorous field experiment on human-AI collaboration in knowledge work. The study was preregistered, conducted with 758 consultants at Boston Consulting Group --- not a laboratory convenience sample, not a blog post, but a controlled experiment with one of the world's premier professional services firms.

The researchers assigned consultants to eighteen realistic consulting tasks and measured the results against a control group working without AI assistance. The findings established two sharply distinct performance regimes.

\begin{evidencebox}[title={Inside the Frontier}]
For the eighteen tasks that fell within AI capability:
\begin{itemize}
  \item Consultants using GPT-4 completed \textbf{12.2 per cent more tasks}
  \item They finished \textbf{25.1 per cent faster}
  \item They produced \textbf{40 per cent higher quality} output, as rated by blind evaluators
\end{itemize}
\end{evidencebox}

\begin{cautionbox}[title={Outside the Frontier}]
For a complex business judgement task requiring integration of real-world context:
\begin{itemize}
  \item AI users were \textbf{19 per cent less likely} to produce correct solutions than those working without AI
  \item The frontier was unpredictable: adjacent tasks that appeared similar fell on opposite sides
\end{itemize}
\end{cautionbox}

The metaphor of a ``jagged'' frontier is precise. Unlike a smooth gradient where tasks become progressively harder for AI, the boundary is irregular and counterintuitive. Two tasks that appear nearly identical from the outside --- similar in domain, complexity, and required expertise --- can fall on opposite sides of the frontier. One is dramatically improved by AI assistance; the other is degraded by it. There is no reliable way to determine which is which without empirical testing.

\begin{figure}[h]
\centering
\fbox{\parbox{0.8\textwidth}{\centering\vspace{3cm}\textbf{[FIGURE: The Jagged Frontier --- Task Performance Map]}\\\textit{Illustrating the unpredictable boundary between AI-capable and AI-incapable tasks}\vspace{3cm}}}
\caption{The Jagged Technological Frontier}
\label{fig:jagged-frontier}
\end{figure}

\section{The Skill-Levelling Effect}

The BCG experiment revealed a second finding with profound implications for organisational design. AI did not improve all consultants equally:

\begin{itemize}
  \item Bottom-half performers saw a \textbf{43 per cent performance jump} with AI assistance.
  \item Top performers improved measurably but by a smaller margin.
  \item AI compressed the talent distribution --- bringing lower-performing consultants much closer to the performance of their higher-performing colleagues.
\end{itemize}

This skill-levelling effect changes what managers are managing. In a pre-AI organisation, individual performance variation is a primary concern: identifying, developing, and retaining high performers is a core management function. When AI compresses the performance distribution, the variation between individuals shrinks. The manager's value shifts from talent differentiation to system coherence --- ensuring that the newly capable workforce is pointed at the right problems and that AI assistance is applied where it helps rather than where it harms.

\section{Three Collaboration Modes}

A follow-up study by \textcite{Randazzo2024}, conducted with 244 BCG consultants, identified three distinct modes of human-AI collaboration. These modes are not merely descriptive; they are predictive of performance outcomes.

\begin{table}[h]
\centering
\caption{Human-AI collaboration modes and performance outcomes}
\label{tab:collaboration-modes}
\begin{tabularx}{\textwidth}{l L{5cm} X}
\toprule
\textbf{Mode} & \textbf{Description} & \textbf{Performance} \\
\midrule
Centaurs & Clear division of labour. The human performs tasks where human judgment is strongest; the AI handles tasks inside the frontier. Distinct handoff points between human and machine. & High \\
\addlinespace
Cyborgs & Continuous interleaving. Human and AI work on the same analysis, passing work back and forth across the frontier boundary. No clean separation of responsibilities. & High \\
\addlinespace
Self-Automators & Wholesale delegation to AI with minimal oversight. The human provides the initial prompt and accepts the output with little or no review. & Lowest \\
\bottomrule
\end{tabularx}
\end{table}

Both Centaurs and Cyborgs outperformed solo humans and solo AI. Self-Automators produced the weakest results --- precisely because wholesale delegation triggers the vigilance problem examined in the next chapter.

\begin{keyinsight}[title={The Organisational Design Implication}]
This taxonomy maps directly onto the three organisational models examined in Chapter~\ref{ch:three-models}:
\begin{itemize}
  \item Block's intelligence layer model creates \textbf{Self-Automator organisations} --- delegating coordination to an AI world model.
  \item Every's agent colony model creates \textbf{Cyborg organisations} --- continuous interleaving between human and personal agent.
  \item The Dynamic Agentic Mesh creates \textbf{Centaur organisations} --- clear boundaries with judgment brokers at the handoff points.
\end{itemize}
The evidence shows that Self-Automators perform worst. This is the empirical case against wholesale delegation.
\end{keyinsight}

\section{Why the Frontier Matters for Governance}

The jagged frontier has four implications that recur throughout this report.

\textbf{First, it invalidates wholesale automation.} If you cannot predict in advance which tasks AI will improve and which it will degrade, a strategy of replacing human judgment with AI judgment is not bold --- it is uninformed. The 19 per cent performance degradation on tasks outside the frontier is not a temporary limitation that will be resolved by the next model generation. It is a structural feature of deploying any system at the boundary of its capability.

\textbf{Second, it defeats one-time change playbooks.} The frontier moves as models improve. A task that falls outside the frontier with GPT-4 may fall inside it with the next generation. Conversely, increased model capability can create new failure modes at different points on the boundary. A change management plan that maps the frontier today will be outdated within months. This demands continuous adaptation, not phased transformation.

\textbf{Third, it creates a new form of expertise.} The person who knows where the frontier falls in their specific domain --- who has mapped which tasks are reliably improved by AI and which are reliably degraded --- holds knowledge that is genuinely scarce and strategically valuable. This is the \meshterm{judgment broker} role developed in Chapter~\ref{ch:role-transformation}. It is not a euphemism for a diminished middle management function. It is a description of expertise that did not exist before the frontier existed.

\textbf{Fourth, it demands active engagement, not passive oversight.} The standard response to AI risk --- ``keep a human in the loop'' --- assumes that the human will maintain vigilance. The next chapter demonstrates why that assumption fails, and what it means for organisational design.

\section{The Frontier Is Dynamic}

It is important to note a limitation of the evidence. The BCG experiment used GPT-4 --- now two model generations behind current frontier models. The specific task boundaries identified in 2023 have almost certainly shifted. Some tasks that fell outside the frontier are now inside it; others may have developed new failure modes as model capabilities have changed.

However, the \emph{principle} holds regardless of where the specific boundary falls at any given moment. The frontier will always be jagged because AI capability will always be unevenly distributed across the task space of any complex domain. The implication --- that organisations need people who actively probe the boundary rather than assuming they know where it is --- becomes more important, not less, as models improve and the frontier moves.

The publication of the BCG study in \emph{Organization Science} in March 2026 \parencite{DellAcqua2026} provides the peer-reviewed validation that the working paper lacked. The jagged frontier is no longer a provocative hypothesis. It is an established finding in the organisational science literature.
